% answers_short.tex —— 2008 年·普通高等学校招生考试·山东卷·理科 · 简短答案速查
% 由 tools/make_short_answers.py 自 answers.tex 的 \gksolhead/\gkanswer 提取，
% 并与 2008/2008_数学真题_山东理_答案要点.md 逐题比对后生成，勿手改（改源 MD 后重跑）。

\begin{gknotice}
\textbf{使用说明}：本节只给最终答案便于快速核对，完整推导见解析版。答案为编写者独立推导，\textbf{非官方标准答案}。
\end{gknotice}

\gkpart{选择题}
\begin{center}\begin{tabular}{ccccccccccccc}
\toprule
题号 & 1 & 2 & 3 & 4 & 5 & 6 & 7 & 8 & 9 & 10 & 11 & 12 \\
\midrule
答案 & B & D & A & A & C & C & B & B & C & A & B & C \\
\bottomrule
\end{tabular}\end{center}

\gkpart{填空题}
\begin{itemize}[leftmargin=2.6em]
  \item[\textbf{13.}] $4$。
  \item[\textbf{14.}] $\dfrac{\sqrt3}3$。
  \item[\textbf{15.}] $\dfrac{\pi}6$。
  \item[\textbf{16.}] $(5,7)$。
\end{itemize}

\gkpart{解答题}
\begin{itemize}[leftmargin=2.6em]
  \item[\textbf{17.}] （1）$\sqrt2$；（2）$\left[\dfrac{2\pi}3+4k\pi,\dfrac{8\pi}3+4k\pi\right]$（$k\in\mathbf Z$）。
  \item[\textbf{18.}] （1）分布列见解析，$E\xi=2$；（2）$\dfrac{34}{243}$。
  \item[\textbf{19.}] （1）$\left\{\dfrac1{S_n}\right\}$ 是公差 $\dfrac12$ 的等差数列，$b_1=1$，$b_n=-\dfrac{2}{n(n+1)}$（$n\geqslant2$）；（2）$-\dfrac{2(2^k-1)}{k(k+1)}$。
  \item[\textbf{20.}] （1）证明见解析；（2）$\dfrac{\sqrt{15}}5$。
  \item[\textbf{21.}] （1）$a>0$ 时极小值 $\dfrac a2\left(1+\ln\dfrac2a\right)$，$a\leqslant0$ 时无极值；（2）证明见解析。
  \item[\textbf{22.}] （1）证明见解析；（2）$x^2=2y$ 或 $x^2=4y$；（3）不存在，理由见解析。
\end{itemize}
